\Needspace{14\baselineskip}
\part{Critical escape and the response generated by the equation}

Part~I has reduced a proposed finite maximal time to one exact terminal
problem.  The kinetic energy remains bounded, while the scale-invariant
\(\dot H^{1/2}\) size must make arbitrarily large excursions arbitrarily
close to the endpoint.  Yet the proof still has to control every required
finite derivative, make the different finite views approach one common
state, and preserve the differentiated velocity--pressure laws there.  The
question is how one fixed evolution could drive a critical quantity through
unbounded levels in the time still available.  We first make that escape
visible in space and time.  Its rapid variation forces the Eulerian response
generated by the original equation, and that response opens the two
independent derivative directions that the terminal estimate must control.

\section{The critical height at a finite maximal time}

The three-dimensional Navier--Stokes scaling is
\begin{equation*}
 u_\lambda(t,x)=\lambda u(\lambda^2t,\lambda x),
 \qquad
 p_\lambda(t,x)=\lambda^2p(\lambda^2t,\lambda x).
\end{equation*}
It preserves the equation and the divergence constraint.  The time interval
is shortened by the factor \(\lambda^{-2}\), and the velocity is amplified by
the factor \(\lambda\).  A spatial norm left unchanged by this transformation
measures the solution at the same scale as the equation.

Let \(\Lambda=(-\Delta)^{1/2}\).  The rescaled Fourier transform is
\begin{equation*}
 \widehat{u_\lambda}(t,\xi)
 =\lambda^{-2}\widehat u(\lambda^2t,\xi/\lambda).
\end{equation*}
Changing variables \(\xi=\lambda\eta\) gives, for every real \(s\) for which
the homogeneous norm is finite,
\begin{align*}
 \norm{u_\lambda(t)}_{\dot H^s}^2
 &=\int_{\R^3}|\xi|^{2s}\lambda^{-4}
   \left|\widehat u(\lambda^2t,\xi/\lambda)\right|^2\,d\xi
 \\
 &=\lambda^{2s-1}
   \int_{\R^3}|\eta|^{2s}
   \left|\widehat u(\lambda^2t,\eta)\right|^2\,d\eta
 \\
 &=\lambda^{2s-1}
   \norm{u(\lambda^2t)}_{\dot H^s}^2.
\end{align*}
The scaling exponent vanishes at \(s=1/2\).  We accordingly define the
critical height
\begin{equation*}
 H(t):=\frac12\norm{\Lambda^{1/2}u(t)}_2^2.
\end{equation*}
The factor \(1/2\) will remove an inessential factor of two when the height is
differentiated.

Scaling explains why this height belongs at the critical frontier.  The
all-data endpoint theorem recalled in Part~I makes it the exact quantity to
follow on a finite maximal branch.  The critical Sobolev embedding gives
\begin{equation*}
 \norm{u(t)}_3
 \le C_{\mathrm{Sob}}\norm{\Lambda^{1/2}u(t)}_2
 =C_{\mathrm{Sob}}\sqrt{2H(t)}.
\end{equation*}
If \(H\) remained bounded near \(T_*\), then the \(L^3\) norm would remain
bounded there as well, contrary to the Escauriaza--Seregin--\v{S}ver\'ak
condition \eqref{eq:ESS-blowup}.  Hence, for the full Schwartz datum class,
\begin{equation}
 T_*<\infty
 \quad\Longrightarrow\quad
 \limsup_{t\uparrow T_*}H(t)=\infty.
 \label{eq:critical-height-limsup}
\end{equation}
Equivalently, for every finite level \(L\) and every
\(\tau<T_*\), there is a time \(t\in(\tau,T_*)\) at which \(H(t)>L\).
The statement allows the height to return below \(L\) between successive
excursions.

For the smooth compactly supported data covered by Seregin's stronger
terminal theorem, Part~I records the full limit
\eqref{eq:seregin-critical-divergence}.  In that historical subcase, every
fixed level is eventually left behind.  The all-Schwartz argument below rests
entirely on the weaker conclusion: every crossing is selected from
\eqref{eq:critical-height-limsup}.

The same height remains tied to the kinetic-energy identity by a direct
Fourier interpolation.  Plancherel and Cauchy--Schwarz give
\begin{align*}
 2H(t)
 &=\int_{\R^3}|\xi|\,|\widehat u(t,\xi)|^2\,d\xi
 \\
 &\le
 \left(\int_{\R^3}|\widehat u(t,\xi)|^2\,d\xi\right)^{1/2}
 \left(\int_{\R^3}|\xi|^2|\widehat u(t,\xi)|^2\,d\xi\right)^{1/2}
 \\
 &=\norm{u(t)}_2\norm{\nabla u(t)}_2.
\end{align*}
After squaring and integrating over \((0,T)\), the energy identity
\eqref{eq:energy} yields
\begin{align*}
 \int_0^T H(t)^2\,dt
 &\le \frac14
 \sup_{0\le t\le T}\norm{u(t)}_2^2
 \int_0^T\norm{\nabla u(t)}_2^2\,dt
 \\
 &\le \frac{\norm{u_0}_2^4}{8\nu}.
\end{align*}
Monotone convergence in \(T\) gives the whole-branch budget
\begin{equation}
 \int_0^{T_*}H(t)^2\,dt
 \le \frac{\norm{u_0}_2^4}{8\nu}<\infty.
 \label{eq:critical-height-budget}
\end{equation}
Thus the same function satisfies two sharply different requirements on a
finite singular branch:
\begin{equation*}
 H\in L^2(0,T_*),
 \qquad
 \limsup_{t\uparrow T_*}H(t)=\infty.
\end{equation*}
The first limits the aggregate time available at a large height.  The second
forces arbitrarily late excursions through every large value before the
finite time \(T_*\).

\subsection{Critical height as a frequency moment}

The difference between kinetic energy and critical height is especially clear
in a dyadic decomposition.  Choose smooth Littlewood--Paley multipliers
\(P_K\), indexed by \(K=2^j\), whose symbols have uniformly finite overlap and
are supported where \(K/2\le |\xi|\le2K\)\bibref{BCD}.  Put
\begin{equation*}
 A_K(t):=\frac12\norm{P_Ku(t)}_2^2.
\end{equation*}
On the support of \(P_K\), the multiplier \(|\xi|^s\) is comparable to
\(K^s\).  Plancherel gives
\begin{equation*}
 \frac12\norm{\Lambda^sP_Ku(t)}_2^2
 \simeq K^{2s}A_K(t),
\end{equation*}
with constants independent of \(K\).  Summing over the finite overlap gives
\begin{equation*}
 \frac12\norm{u(t)}_2^2\simeq\sum_KA_K(t),
 \qquad
 H(t)\simeq\sum_KK A_K(t).
\end{equation*}
Kinetic energy is the unweighted shell mass.  Critical height is its first
frequency moment.  A bounded unweighted mass and an unbounded first moment
force nonzero contributions at arbitrarily large frequencies: for any fixed
\(K_0\),
\begin{equation*}
 \sum_{K\le K_0}K A_K(t)
 \le K_0\sum_KA_K(t)
 \lesssim K_0\norm{u_0}_2^2,
\end{equation*}
so the portion below \(K_0\) remains bounded uniformly in time.  Whenever
\(H(t)\) exceeds that bound, frequencies above \(K_0\) must contribute.

This calculation fixes the spatial burden: a bounded amount of kinetic energy
must contribute at frequencies beyond every prescribed cutoff.  A finite
terminal time creates a second burden, because those high-frequency
excursions must occur before the remaining time is exhausted.  We follow that
temporal compression first.  The shell weights return after the
critical-height balance has identified the response that supplies the rise.

\subsection{Finite-time compression}

Estimate~\eqref{eq:critical-height-budget} immediately controls the total time
spent above a large height.  For every \(L>0\), Chebyshev's inequality gives
\begin{equation}
 \bigl|\{t\in(0,T_*):H(t)\ge L\}\bigr|
 \le
 \frac{\norm{u_0}_2^4}{8\nu L^2}.
 \label{eq:height-occupancy}
\end{equation}
The terminal limsup and the occupancy estimate
\eqref{eq:height-occupancy} perform different work.  The
limsup supplies arbitrarily late excursions above every high level and permits
returns below it.  The occupancy estimate bounds the aggregate time available
above that level.  We will not need the height to stay above a level once it
has crossed it.  First-passage endpoints alone will give a fixed gain across
intervals whose lengths tend to zero.

The crossing times follow directly from continuity on the classical interval.

\begin{lemma}[First-passage compression]
\label{lem:first-passage-compression}
Assume \(T_*<\infty\).  Let \(L_j\uparrow\infty\) be strictly increasing with
\(L_1>H(0)\), and define
\begin{equation*}
 t_j:=\inf\{t\in(0,T_*):H(t)=L_j\}.
\end{equation*}
Then every \(t_j\) is well defined,
\begin{equation*}
 t_1<t_2<\cdots<T_*,
 \qquad
 t_j\uparrow T_*.
\end{equation*}
\end{lemma}

\begin{proof}
The height is continuous on every compact subinterval of the classical
lifespan.  The limsup in \eqref{eq:critical-height-limsup} places a time
\(s_j<T_*\) with \(H(s_j)>L_j\).  Since \(H(0)<L_1\le L_j\), the intermediate
value theorem supplies a point in \([0,s_j]\) where \(H=L_j\).  The level set
is closed inside that compact interval, so it has a first point \(t_j\).

Continuity also fixes the order of the crossings.  A path from
\(H(0)<L_j<L_{j+1}\) to the level \(L_{j+1}\) must cross \(L_j\) first, so
\(t_j<t_{j+1}\).  The increasing sequence has a limit \(\tau\le T_*\).  If
\(\tau<T_*\), choose \(\varepsilon>0\) with
\(\tau+\varepsilon<T_*\).  Classicality bounds \(H\) on the compact interval
\([0,\tau+\varepsilon]\), while
\(H(t_j)=L_j\to\infty\) inside that interval.  This is impossible.  Hence
\(\tau=T_*\).
\end{proof}

Take unit-spaced levels,
\begin{equation*}
 L_j=L_1+j-1,
\end{equation*}
and set
\begin{equation*}
 \delta_j:=t_{j+1}-t_j.
\end{equation*}
The positive gaps fit inside the finite interval and satisfy
\(\delta_j\to0\).  On each closed interval
\([t_j,t_{j+1}]\), the solution is classical, so the fundamental theorem of
calculus gives
\begin{align*}
 1
 &=H(t_{j+1})-H(t_j)
 \\
 &=\int_{t_j}^{t_{j+1}}H'(s)\,ds
 \\
 &\le\int_{t_j}^{t_{j+1}}|H'(s)|\,ds.
\end{align*}
The average absolute response obeys
\begin{equation}
 \frac1{\delta_j}
 \int_{t_j}^{t_{j+1}}|H'(s)|\,ds
 \ge\frac1{\delta_j}
 \longrightarrow\infty.
 \label{eq:zeno-average}
\end{equation}
A fixed gain is being carried across intervals of vanishing length.  We call
the divergence in \eqref{eq:zeno-average} the Zeno
compression of the critical height.  Only the endpoint values are used: the
height may fall below \(L_j\) inside the interval and after \(t_j\).

The derivative in this calculation is not an abstract scalar source.  Since
\[
 H(t)=\frac12\norm{\Lambda^{1/2}u(t)}_2^2,
\]
classical differentiation gives
\begin{equation*}
 H'(t)=\operatorname{Re}
 \pair{\Lambda^{1/2}u_t(t)}{\Lambda^{1/2}u(t)}.
\end{equation*}
The compressed rise is carried by the first Eulerian response \(u_t\) acting
against the same velocity field whose critical size is rising.  Physical
acceleration, however, is the material derivative \(D_tu\), not \(u_t\).
Understanding the response requires us to separate those two derivatives
without discarding the transport correction between them.

The unit ladder measures temporal compression.  Repeating the first-passage
construction with a dyadic ladder will measure the work performed by the
equation.  If
\(L_{j+1}=2L_j\), then
\begin{equation*}
 H(t_{j+1})-H(t_j)=L_j.
\end{equation*}
After the critical-height balance has been derived, integrating it over these
same intervals will show that the nonlinear pressure-completed transfer must
supply this entire gain together with the intervening viscous loss.

The terminal threat has now become a response problem.  The height records the
size of the event, while the equation must identify the fields that produce
its change.  The physical route is velocity, acceleration, and jerk.  The
analytic route must express that hierarchy at fixed spatial coordinates,
where the whole-field Sobolev norms are measured.

\section{From acceleration and jerk to the Eulerian tower}

Part~I fixed the material derivative along particle trajectories and showed
that finite-time escape must be accumulated by higher material response.  The
critical height, however, is a norm of the entire field at one fixed time.  We
now compute how material acceleration and jerk decompose into the Eulerian rows
\(V_n=\partial_t^n u\).

The next material response is the jerk
\begin{equation*}
 D_t^2u=D_t(D_tu).
\end{equation*}
It is already more involved than a second partial derivative in time.  Set
\begin{equation*}
 V_n:=\partial_t^nu,
 \qquad n\in\Nzero.
\end{equation*}
Since
\begin{equation*}
 D_tu=V_1+(V_0\cdot\nabla)V_0,
\end{equation*}
one more application of \(D_t\) gives
\begin{align*}
 D_t^2u
 &=(\partial_t+V_0\cdot\nabla)
   \bigl(V_1+(V_0\cdot\nabla)V_0\bigr)
 \\
 &=V_2
   +(V_1\cdot\nabla)V_0
   +(V_0\cdot\nabla)V_1
   +(V_0\cdot\nabla)V_1
   +(V_0\cdot\nabla)\bigl((V_0\cdot\nabla)V_0\bigr)
 \\
 &=V_2
   +(V_1\cdot\nabla)V_0
   +2(V_0\cdot\nabla)V_1
   +(V_0\cdot\nabla)\bigl((V_0\cdot\nabla)V_0\bigr).
\end{align*}
The row \(V_2\) is one part of the material jerk.  The other terms are the
exact transport corrections created by changing from a moving particle to a
fixed spatial point.  At higher order, this difference keeps growing through
the commutator of the material derivative with spatial differentiation.  For
example,
\begin{equation*}
 [D_t,\partial_\ell]f
 =-(\partial_\ell u\cdot\nabla)f.
\end{equation*}
This commutator explains why the proof does not build its mixed array with
successive powers of \(D_t\).  Material differentiation continually creates
new spatial correction terms.  By contrast, \(\partial_t\) commutes with every
fixed spatial derivative, while the differentiated transport sum records all
of the corrections explicitly.  The Eulerian rows keep temporal depth and
spatial depth as separate coordinates without losing the physics that related
them.

The acceleration and jerk calculation separates the moving-particle response
from the fixed-coordinate derivatives used by the proof.  The critical height
is a norm of the whole field at one fixed time, and its Fourier multipliers act
in fixed spatial coordinates.  Global norm divergence leaves every
particle-path choice unresolved.  We thus work with the Eulerian rows
\(V_n\), while the differentiated transport terms retain exactly the
corrections created by the change of frame.

\subsection{The exact escalation from one Eulerian row to the next}

The fundamental theorem of calculus turns the response intuition into a norm
statement.  Let \(X\) be a Banach space, fix
\(0\le r<t<T_*\), and suppose
\begin{equation*}
 V_j\in W^{1,1}(r,t;X),
 \qquad
 \partial_tV_j=V_{j+1}.
\end{equation*}
The Bochner fundamental theorem gives
\begin{align}
 V_j(t)-V_j(r)
 &=\int_r^tV_{j+1}(s)\,ds,
 \notag\\
 \norm{V_j(t)}_X
 &\le
 \norm{V_j(r)}_X
 +\int_r^t\norm{V_{j+1}(s)}_X\,ds.
 \label{eq:temporal-escalation}
\end{align}
If \(V_{j+1}\in L^1(0,T_*;X)\), the right side stays bounded uniformly in
\(t<T_*\).  On the finite interval, Cauchy--Schwarz gives
\begin{equation*}
 \int_0^{T_*}\norm{V_{j+1}(s)}_X\,ds
 \le
 \sqrt{T_*}\,\norm{V_{j+1}}_{L^2(0,T_*;X)}.
\end{equation*}
Thus one whole-terminal \(L^2\) row supplies finite total variation to the row
below it.  The implication is conditional at each level: if a chosen row
escapes, the next row cannot have finite accumulated variation in the same
space.  Escape of \(V_0\) alone does not prove that every higher \(V_j\)
diverges.  The equation generates all of those rows for a different reason:
smooth continuation must retain every finite time derivative at the endpoint.

Apply this observation at the critical spatial order.  Equation
\eqref{eq:critical-height-limsup} is equivalent to
\begin{equation*}
 \limsup_{t\uparrow T_*}
 \norm{V_0(t)}_{\dot H^{1/2}}=\infty.
\end{equation*}
If the hypotheses needed for \eqref{eq:temporal-escalation} hold with
\(X=\dot H^{1/2}\), then
\begin{equation*}
 \int_0^{T_*}\norm{V_1(s)}_{\dot H^{1/2}}\,ds=\infty.
\end{equation*}
On every strict classical subinterval the same argument applies to any fixed
row for which an escape has been established.  At the base row it identifies
the exact channel through which critical escape must pass.  Whole-terminal
control of \(V_1\) in a sufficiently strong space would close that channel;
retaining every higher row will later give the complete endpoint jet needed
for a smooth restart.

\subsection{The pressure-complete temporal recurrence}

The Eulerian hierarchy becomes an object of the equation once pressure is
differentiated with it.  Define
\begin{equation*}
 Q_n:=\partial_t^np.
\end{equation*}
The zeroth temporal row is the original momentum equation:
\begin{equation*}
 V_1+(V_0\cdot\nabla)V_0+\nabla Q_0
 =\nu\Delta V_0.
\end{equation*}
Differentiating once gives
\begin{equation*}
 V_2
 +(V_1\cdot\nabla)V_0
 +(V_0\cdot\nabla)V_1
 +\nabla Q_1
 =\nu\Delta V_1.
\end{equation*}
The two transport terms have different origins.  In the first, the derivative
lands on the velocity doing the carrying.  In the second, it lands on the
velocity being carried.  Both remain in the same differentiated row.

At order \(n\), Leibniz' rule distributes the time derivatives across the two
velocity factors:
\begin{align*}
 \partial_t^n\bigl((u\cdot\nabla)u\bigr)
 &=\sum_{a=0}^n\binom na
   (\partial_t^au\cdot\nabla)\partial_t^{n-a}u
 \\
 &=\sum_{a=0}^n\binom na
   (V_a\cdot\nabla)V_{n-a}.
\end{align*}
The coefficient \(\binom na\) counts the choices of the \(a\) derivatives
landing on the first velocity factor.

Pressure at the same order is fixed by incompressibility.  Taking divergence
of the differentiated momentum row and using
\(\nabla\cdot V_a=0\) gives
\begin{align*}
 -\Delta Q_n
 &=\nabla\cdot
   \sum_{a=0}^n\binom na(V_a\cdot\nabla)V_{n-a}
 \\
 &=\partial_i\partial_j
   \sum_{a=0}^n\binom naV_{a,i}V_{n-a,j}.
\end{align*}
With the decaying normalization fixed in \eqref{eq:pressure}, inversion by
Riesz transforms gives
\begin{equation}
 Q_n
 =R_iR_j\sum_{a=0}^n\binom naV_{a,i}V_{n-a,j}.
 \label{eq:pressure-time-jet}
\end{equation}
Substitution into the momentum row produces the temporal recurrence
\begin{equation}
 V_{n+1}
 =\nu\Delta V_n
 -\sum_{a=0}^n\binom na(V_a\cdot\nabla)V_{n-a}
 -\nabla Q_n,
 \qquad
 \nabla\cdot V_n=0.
 \label{eq:velocity-time-jet}
\end{equation}

Equations~\eqref{eq:velocity-time-jet} and
\eqref{eq:pressure-time-jet} preserve the same VPI response at temporal depth
\(n\).  Viscosity
advances the spatial curvature through \(\Delta V_n\).  Every lower temporal
split appears in the transport sum.  The pressure equation performs the
whole-space correction required to keep \(V_{n+1}\) divergence-free.  These
three operations generate one next row.

Pressure acts simultaneously across the whole field.  At the base row, the
Riesz formula has the Calder\'on--Zygmund representation
\begin{equation*}
 p(t,x)=-\frac13|u(t,x)|^2
 +\operatorname{p.v.}\frac1{4\pi}\int_{\R^3}
 \frac{3z_i z_j-|z|^2\delta_{ij}}{|z|^5}\,
 u_i(t,y)u_j(t,y)\,dy,
 \qquad z=x-y.
\end{equation*}
The local term is the delta contribution in the distributional Hessian of the
Newtonian kernel; together with the principal-value term it is exactly
\eqref{eq:pressure}.  A change in the velocity far from \(x\) can still change
the pressure gradient at \(x\), and the differentiated Riesz formula carries
that same global coupling through every \(Q_n\).  Because the Riesz transforms
have order zero on Sobolev spaces, the upper velocity rows of a finite
rectangle will also determine the size and normalization of every pressure row
it contains; no independent pressure datum will be added at the endpoint.

The recurrence also requires separate temporal and spatial indices.  The
Laplacian differentiates \(V_n\) twice in space, transport
contains a spatial derivative of one factor, and pressure is obtained from an
order-zero nonlocal operator applied to all lower products.  Each temporal row
opens an entire spatial column.  Before constructing the full two-index array,
we return to the critical height.  The scalar height balance will cancel a pure
pressure gradient, but its differentiated interaction terms are still computed
from the pressure-complete rows above.  This lets the height read spatial
information from the response without replacing the response itself.

\section{The completed critical-height rows}

For \(s\ge0\), define the homogeneous Sobolev energy
\begin{equation*}
 E_s(t):=\frac12\norm{\Lambda^su(t)}_2^2.
\end{equation*}
On every compact classical subinterval,
\begin{equation*}
 E_s'(t)
 =\operatorname{Re}
  \pair{\Lambda^su(t)}{\Lambda^su_t(t)}.
\end{equation*}
Insert the momentum equation.  The viscous term gives
\begin{align*}
 \nu\operatorname{Re}
 \pair{\Lambda^su}{\Lambda^s\Delta u}
 &=-\nu\norm{\Lambda^{s+1}u}_2^2
 \\
 &=-2\nu E_{s+1}.
\end{align*}
The pressure term cancels in this whole-space pairing because the velocity is
divergence-free:
\begin{align*}
 \operatorname{Re}
 \pair{\Lambda^su}{\Lambda^s\nabla p}
 &=\operatorname{Re}
 \pair{\Lambda^{2s}u}{\nabla p}
 \\
 &=-\operatorname{Re}
 \pair{\nabla\cdot\Lambda^{2s}u}{p}
 \\
 &=0.
\end{align*}
The cancellation is a consequence of the pressure equation and
incompressibility.  It removes the pure gradient from this one scalar pairing;
it does not remove pressure from the differentiated field.  Time derivatives
of the interaction below contain the rows \(V_a\) and the pressure responses
\(Q_a\) fixed by \eqref{eq:pressure-time-jet}--\eqref{eq:velocity-time-jet}.
In the Sobolev balance, the VPI law separates into the viscous term above and
the signed pressure-complete interaction current
\begin{equation*}
 \mathcal P_s
 :=-\operatorname{Re}
 \pair{\Lambda^su}
 {\Lambda^s\bigl((u\cdot\nabla)u+\nabla p\bigr)}.
\end{equation*}
Equivalently,
\begin{equation*}
 \mathcal P_s
 =-\operatorname{Re}
 \pair{\Lambda^su}
 {\Lambda^s\PP\bigl((u\cdot\nabla)u\bigr)},
\end{equation*}
where \(\PP\) is the Leray projection.  The pressure-explicit and projected
forms are two global pairings of the same incompressible interaction.  The
Sobolev balance is
\begin{equation}
 E_s'=-2\nu E_{s+1}+\mathcal P_s.
 \label{eq:Es-balance}
\end{equation}

At \(s=0\), the transport contribution also cancels and
\eqref{eq:Es-balance} is the kinetic-energy identity.  At the critical order
\(s=1/2\), it becomes
\begin{equation}
 H'=-2\nu E_{3/2}+\mathcal P_{1/2}.
 \label{eq:H1}
\end{equation}
The two terms on the right have clear roles.  Viscosity removes critical
height through the next spatial energy \(E_{3/2}\).  The signed interaction
\(\mathcal P_{1/2}\) supplies whatever height gain remains after that loss.

\subsection{The first-passage intervals enter the equation}

Integrating \eqref{eq:H1} over an arbitrary classical interval \([r,t]\)
gives
\begin{equation}
 H(t)-H(r)
 +2\nu\int_r^tE_{3/2}(s)\,ds
 =\int_r^t\mathcal P_{1/2}(s)\,ds.
 \label{eq:critical-passage-balance}
\end{equation}
Return first to the unit-spaced crossings from
Lemma~\ref{lem:first-passage-compression}.  Since
\(H(t_{j+1})-H(t_j)=1\), the exact balance on the shrinking interval is
\begin{equation}
 \int_{t_j}^{t_{j+1}}\mathcal P_{1/2}(s)\,ds
 =1+2\nu\int_{t_j}^{t_{j+1}}E_{3/2}(s)\,ds.
 \label{eq:unit-passage-work}
\end{equation}
The transfer supplies the unit height gain and every unit of critical
dissipation incurred while that gain is made.  The intervals have lengths
\(\delta_j\to0\), so the same fixed gain and dissipation are being supplied on
a sequence of vanishing time scales.

The dyadic levels separate the amplitude of that work from the shrinking
interval.  If \(L_{j+1}=2L_j\), then
\(H(t_{j+1})-H(t_j)=L_j\), and
\eqref{eq:critical-passage-balance} gives
\begin{equation}
 \int_{t_j}^{t_{j+1}}\mathcal P_{1/2}(s)\,ds
 =L_j+2\nu\int_{t_j}^{t_{j+1}}E_{3/2}(s)\,ds
 \ge L_j.
 \label{eq:dyadic-passage-work}
\end{equation}
The integrated transfer must grow with the level being crossed.
Equation~\eqref{eq:dyadic-passage-work} has converted the shrinking-time
picture into an exact demand on the equation: over the interval that raises
the height from \(L_j\) to \(2L_j\), the signed interaction must supply at
least \(L_j\), in addition to the viscous loss.  The identity determines that
demand; an upper estimate must still control it.

At critical order, the response has the same spatial scale as the viscous
term.  Self-adjointness of \(\Lambda^{1/2}\), the pressure cancellation,
H\"older's inequality, and the fractional Sobolev embedding give
\begin{align}
 |\mathcal P_{1/2}|
 &=\left|\bigl((u\cdot\nabla)u,\Lambda u\bigr)_2\right|
 \notag\\
 &\le\norm{u}_3\norm{\nabla u}_3\norm{\Lambda u}_3
 \notag\\
 &\le C\norm{u}_{\dot H^{1/2}}
          \norm{\Lambda^{3/2}u}_2^2.
 \label{eq:critical-response-estimate}
\end{align}
For sufficiently small critical data, this estimate allows the viscous term
to absorb the interaction, which is the familiar Fujita--Kato mechanism
\bibref{FujitaKato,Kato}.  The arbitrary datum in the main theorem lies beyond
that smallness regime.  We cannot close the first balance by estimating its
interaction as a small error.  The next move is to differentiate the complete
identity instead, retaining every response term it creates.  The passage
identities \eqref{eq:unit-passage-work} and
\eqref{eq:dyadic-passage-work} remain valid at every critical size and tell us
what those completed rows must eventually control.

\subsection{One height derivative, two exact descriptions}

The first description is kinematic.  Since
\begin{equation*}
 H(t)=\frac12
 \pair{\Lambda^{1/2}V_0(t)}{\Lambda^{1/2}V_0(t)},
\end{equation*}
the first derivative is
\begin{equation*}
 H'(t)=\operatorname{Re}
 \pair{\Lambda^{1/2}V_1(t)}{\Lambda^{1/2}V_0(t)}.
\end{equation*}
Repeated Leibniz differentiation gives, for every \(n\in\Nzero\),
\begin{equation}
 H^{(n)}(t)
 =\frac12\sum_{a=0}^n\binom na
 \operatorname{Re}
 \pair{\Lambda^{1/2}V_a(t)}
 {\Lambda^{1/2}V_{n-a}(t)}.
 \label{eq:height-temporal-contraction}
\end{equation}
The \(n\)-th derivative of the scalar height contains the complete finite
temporal block
\begin{equation*}
 V_0,V_1,\ldots,V_n.
\end{equation*}
It contracts that block to one number.  The block still carries much more
information than the scalar contraction.

The second description is dynamic.  Differentiate
\eqref{eq:Es-balance} and substitute the same balance at each new spatial
order.  The first four rows are
\begin{align*}
 H
 &=E_{1/2},
 \\
 H'
 &=-2\nu E_{3/2}+\mathcal P_{1/2},
 \\
 H''
 &=4\nu^2E_{5/2}
   -2\nu\mathcal P_{3/2}
   +\partial_t\mathcal P_{1/2},
 \\
 H'''
 &=-8\nu^3E_{7/2}
   +4\nu^2\mathcal P_{5/2}
   -2\nu\partial_t\mathcal P_{3/2}
   +\partial_t^2\mathcal P_{1/2}.
\end{align*}
Each differentiation advances the leading viscous energy by one full spatial
derivative.  It also differentiates every interaction current already present
and creates the current belonging to the new spatial level.  The current terms
are part of the identity at the same status as the leading energy.

\begin{lemma}[Completed critical-height rows]
\label{lem:height-tower}
For every \(n\in\Nzero\),
\begin{equation}
 H^{(n)}
 =(-2\nu)^nE_{n+1/2}
 +\sum_{j=0}^{n-1}(-2\nu)^j
   \partial_t^{\,n-1-j}\mathcal P_{j+1/2},
 \label{eq:height-tower}
\end{equation}
where the sum is empty for \(n=0\).  Define the completed row
\begin{equation}
 \mathscr C_n
 :=\frac{
 H^{(n)}
 -\displaystyle\sum_{j=0}^{n-1}(-2\nu)^j
   \partial_t^{\,n-1-j}\mathcal P_{j+1/2}}
 {(-2\nu)^n}.
 \label{eq:completed-row}
\end{equation}
Then
\begin{equation*}
 \mathscr C_n=E_{n+1/2}.
\end{equation*}
\end{lemma}

\begin{proof}
For \(n=0\), the assertion is the definition
\(H=E_{1/2}\).  Suppose \eqref{eq:height-tower} holds at order \(n\).
Differentiate its leading term and apply \eqref{eq:Es-balance} at
\(s=n+1/2\):
\begin{align*}
 (-2\nu)^nE_{n+1/2}'
 &=(-2\nu)^n
   \bigl(-2\nu E_{n+3/2}+\mathcal P_{n+1/2}\bigr)
 \\
 &=(-2\nu)^{n+1}E_{n+3/2}
   +(-2\nu)^n\mathcal P_{n+1/2}.
\end{align*}
The derivative of the existing response sum is
\begin{equation*}
 \sum_{j=0}^{n-1}(-2\nu)^j
 \partial_t^{\,n-j}\mathcal P_{j+1/2}.
\end{equation*}
The newly created term
\((-2\nu)^n\mathcal P_{n+1/2}\) is exactly the missing \(j=n\) term.  The
sum now runs from \(j=0\) through \(j=n\), and the leading energy is
\((-2\nu)^{n+1}E_{n+3/2}\).  This is
\eqref{eq:height-tower} at order \(n+1\).  Solving the identity for its
leading energy gives \eqref{eq:completed-row}.
\end{proof}

The first completed rows read
\begin{equation*}
 \mathscr C_0=E_{1/2},
 \qquad
 \mathscr C_1=E_{3/2},
 \qquad
 \mathscr C_2=E_{5/2},
 \qquad
 \mathscr C_3=E_{7/2}.
\end{equation*}
The completion is mathematically essential.  At the second derivative,
\begin{equation*}
 E_{5/2}
 =\frac{H''+2\nu\mathcal P_{3/2}
              -\partial_t\mathcal P_{1/2}}
        {4\nu^2}.
\end{equation*}
The spatial energy is isolated only after the full differentiated response has
been retained.

The dyadic shell weights from the opening section now return.  Since
\begin{equation*}
 E_{n+1/2}(P_Ku)
 =\frac12\norm{\Lambda^{n+1/2}P_Ku}_2^2,
\end{equation*}
the Littlewood--Paley calculation gives
\begin{equation}
 E_{n+1/2}(u)
 \simeq\sum_KK^{2n+1}A_K.
 \label{eq:completed-shell-weight}
\end{equation}
Equation~\eqref{eq:completed-shell-weight} shows that critical height uses the
weight \(K\).  The next completed row uses \(K^3\),
then \(K^5\), and each increase in temporal depth adds the factor \(K^2\).
The high-frequency contributions required by critical escape are being read
at progressively stronger spatial weights.

\section{The reciprocal temporal--spatial structure}

We can now place the two descriptions of \(H^{(n)}\) beside each other.  The
kinematic identity \eqref{eq:height-temporal-contraction} says that the
\(n\)-th height derivative is built from the temporal block
\[
 V_0,V_1,\ldots,V_n.
\]
The dynamic identity \eqref{eq:height-tower} says that, once every
differentiated interaction current is kept, that same scalar derivative has
the leading spatial energy
\[
 E_{n+1/2}=\frac12\norm{\Lambda^{n+1/2}u}_2^2.
\]
One index is now doing two visible jobs.  It records how far the response has
been differentiated in time, and it records the stronger spatial weight
exposed in the completed identity.

This is the reciprocal structure.  A proposed finite-time singularity would
have to lose spatial control as its active scales contract.  The temporal
response forced by that escape points in the opposite direction: each new
height derivative exposes the next stronger spatial energy of the base
velocity.  The two movements are balanced because they occur inside one exact
identity, and opposing because the singular picture needs spatial regularity
to disappear while the differentiated response keeps revealing higher spatial
orders that a successful estimate could control.

Nothing in the identity has yet supplied that control.  Differentiation does
not create regularity, and the completed-row formula is not an estimate.  It
tells us which spatial energy is hidden inside each temporal response after
the complete VPI current has been retained.  Part~III will later provide the
whole-terminal bounds; the present calculation identifies the structure those
bounds must reach.

There is also an important index distinction.  The formula does not say
\(V_n\in H^{n+1/2}\).  Its leading energy belongs to the base velocity \(u\).
Each temporal row \(V_n\) has its own independently indexed spatial column:
\begin{equation*}
 V_n,\qquad
 \nabla V_n,\qquad
 \nabla^2V_n,\qquad\ldots.
\end{equation*}
The completed height sequence reads one diagonal through the spatial hierarchy
of \(u\); it does not retain the independent temporal and spatial indices.
The scalar contraction also discards the spatial distribution of the field.
To carry smoothness to the endpoint, we need every finite spatial order in
every finite temporal row, with the normalized pressure response attached.

Smoothness asks for every entry to the right of a fixed temporal row.
Simultaneous membership at every finite spatial order gives
\begin{equation*}
 V_n\in\bigcap_{m=0}^\infty H^m(\R^3)=H^\infty(\R^3).
\end{equation*}
For each integer \(k\ge0\), choose
\(m>k+3/2\).  Sobolev embedding gives
\begin{equation*}
 H^m(\R^3)\hookrightarrow C^k(\R^3),
\end{equation*}
so membership at every finite \(m\) gives membership in every \(C^k\):
\begin{equation*}
 H^\infty(\R^3)\hookrightarrow C^\infty(\R^3).
\end{equation*}
The operation is an intersection.  Each finite membership may have its own
constant; smoothness does not require one numerical bound uniform over the
entire infinite array.

The diagonal has shown why spatial depth matters, but it cannot construct the
missing columns.  For that we return to the pressure-complete recurrence
\eqref{eq:pressure-time-jet}--\eqref{eq:velocity-time-jet} and differentiate
each temporal row in space.  This produces the full mixed jet while keeping
the two derivative indices independent.

\section{The pressure-complete mixed jet}

Fix one row of the recurrence
\eqref{eq:pressure-time-jet}--\eqref{eq:velocity-time-jet} and apply every
finite spatial derivative to it.  For a multi-index
\(\alpha=(\alpha_1,\alpha_2,\alpha_3)\), define
\begin{equation*}
 J_{n,\alpha}:=\partial_t^n\partial_x^\alpha u,
 \qquad
 \Pi_{n,\alpha}:=\partial_t^n\partial_x^\alpha p.
\end{equation*}
The temporal index \(n\) counts fixed-point response.  The multi-index
\(\alpha\) counts spatial differentiation.  The two directions remain
independent:
\begin{equation*}
 \begin{array}{c|cccc}
  & |\alpha|=0 & |\alpha|=1 & |\alpha|=2 & \cdots\\
 \hline
  n=0 & u & \nabla u & \nabla^2u & \cdots\\
  n=1 & u_t & \nabla u_t & \nabla^2u_t & \cdots\\
  n=2 & u_{tt} & \nabla u_{tt} & \nabla^2u_{tt} & \cdots\\
  \vdots & \vdots & \vdots & \vdots & \ddots
 \end{array}
\end{equation*}
Moving down changes temporal depth.  Moving right changes spatial depth.  The
completed-height rows describe one diagonal relation inside this array; the
differentiated recurrence generates and retains the array itself.

For \(\beta\le\alpha\), define the multi-index coefficient
\begin{equation*}
 \binom{\alpha}{\beta}
 :=\prod_{\ell=1}^3
   \binom{\alpha_\ell}{\beta_\ell}.
\end{equation*}
Apply \(\partial_x^\alpha\) to one transport term in
\eqref{eq:velocity-time-jet}.  The spatial product rule gives
\begin{align*}
 \partial_x^\alpha
 \bigl((V_a\cdot\nabla)V_{n-a}\bigr)
 &=\sum_{\beta\le\alpha}
   \binom{\alpha}{\beta}
   \bigl(\partial_x^\beta V_a\cdot\nabla\bigr)
   \partial_x^{\alpha-\beta}V_{n-a}
 \\
 &=\sum_{\beta\le\alpha}
   \binom{\alpha}{\beta}
   (J_{a,\beta}\cdot\nabla)
   J_{n-a,\alpha-\beta}.
\end{align*}
Insert this expansion into the temporal binomial sum.  Since spatial
derivatives commute with \(\partial_t\), \(\nabla\), and \(\Delta\), every
mixed row obeys
\begin{equation}
\begin{aligned}
 J_{n+1,\alpha}
 &+\sum_{a=0}^n\sum_{\beta\le\alpha}
 \binom na\binom{\alpha}{\beta}
 (J_{a,\beta}\cdot\nabla)
 J_{n-a,\alpha-\beta}
 +\nabla\Pi_{n,\alpha}
 \\
 &=\nu\Delta J_{n,\alpha},
 \qquad
 \nabla\cdot J_{n,\alpha}=0.
\end{aligned}
\label{eq:mixed-jet}
\end{equation}
The pressure row is fixed by the divergence equation:
\begin{equation}
\begin{aligned}
 -\Delta\Pi_{n,\alpha}
 &=\partial_i\partial_j
   \partial_t^n\partial_x^\alpha(u_i u_j),
 \\
 \Pi_{n,\alpha}
 &=R_iR_j
 \sum_{a=0}^n\sum_{\beta\le\alpha}
 \binom na\binom{\alpha}{\beta}
 J_{a,\beta,i}J_{n-a,\alpha-\beta,j}.
\end{aligned}
\label{eq:mixed-pressure}
\end{equation}
The first binomial factor distributes temporal derivatives.  The second
distributes spatial derivatives.  Every term is produced by the original
quadratic transport and the original pressure normalization.

The first spatial face makes the formula concrete.  Take \(n=0\) and
\(\alpha=e_\ell\).  Equation~\eqref{eq:mixed-jet} becomes
\begin{equation*}
 \partial_t\partial_\ell u
 +(u\cdot\nabla)\partial_\ell u
 +(\partial_\ell u\cdot\nabla)u
 +\nabla\partial_\ell p
 =\nu\Delta\partial_\ell u.
\end{equation*}
The spatial derivative lands once on the transporting velocity and once on
the transported velocity.  At \(n=1\), the same spatial face is
\begin{align*}
 \partial_t^2\partial_\ell u
 &+(u\cdot\nabla)\partial_\ell u_t
  +(\partial_\ell u\cdot\nabla)u_t
 \\
 &+(u_t\cdot\nabla)\partial_\ell u
  +(\partial_\ell u_t\cdot\nabla)u
  +\nabla\partial_\ell p_t
 \\
 &=\nu\Delta\partial_\ell u_t.
\end{align*}
The four transport terms are the four ways one time derivative and one spatial
derivative can be distributed across the two velocity factors.  They remain
differentiated faces of the original solution.

\subsection{The datum generates the initial jet}

The mixed jet begins at the initial datum.  Set
\begin{equation*}
 V_n^0:=V_n(0),
 \qquad
 Q_n^0:=Q_n(0),
 \qquad
 V_0^0=u_0.
\end{equation*}
Evaluating \eqref{eq:pressure-time-jet} and
\eqref{eq:velocity-time-jet} at \(t=0\) gives
\begin{align}
 Q_n^0
 &=R_iR_j\sum_{a=0}^n\binom na
   V_{a,i}^0V_{n-a,j}^0,
 \label{eq:initial-pressure}
 \\
 V_{n+1}^0
 &=\nu\Delta V_n^0
 -\sum_{a=0}^n\binom na
  (V_a^0\cdot\nabla)V_{n-a}^0
 -\nabla Q_n^0.
 \label{eq:initial-velocity}
\end{align}
The first row is already determined entirely by \(u_0\):
\begin{equation*}
 Q_0^0=R_iR_j(u_{0,i}u_{0,j}),
\end{equation*}
and
\begin{equation*}
 V_1^0
 =\nu\Delta u_0
 -(u_0\cdot\nabla)u_0
 -\nabla Q_0^0.
\end{equation*}

The datum class \(\SSigma(\R^3)\) is contained in
\(H^\infty_\sigma(\R^3)\).  Suppose inductively that
\begin{equation*}
 V_0^0,\ldots,V_n^0\in H^\infty.
\end{equation*}
Fix any finite target order \(m\).  Choose an auxiliary Sobolev order high
enough that it is an algebra and remains above \(m\) after every derivative
displayed in \eqref{eq:initial-pressure}--\eqref{eq:initial-velocity}.  The
Sobolev multiplication estimate from
Appendix~\ref{app:analytic-framework}, boundedness of Riesz transforms on
every \(H^m\), and closure under differentiation then give
\begin{equation*}
 Q_n^0,V_{n+1}^0\in H^m.
\end{equation*}
Since \(m\) was arbitrary,
\begin{equation*}
 Q_n^0,V_{n+1}^0\in H^\infty.
\end{equation*}
Induction yields
\begin{equation}
 V_n^0,Q_n^0\in H^\infty(\R^3)
 \qquad(n\in\Nzero).
 \label{eq:initial-all-orders}
\end{equation}
Applying \(\partial_x^\alpha\) gives every
\begin{equation*}
 J_{n,\alpha}(0),
 \qquad
 \Pi_{n,\alpha}(0).
\end{equation*}

Equation~\eqref{eq:initial-all-orders} is the datum-side origin of the mixed
jet.  The recursion determines every initial coordinate from \(u_0\).
Whole-terminal control of those coordinates is a separate conclusion.  A
terminal value for \(V_n\)
requires simultaneous control of \(V_n\) and its time derivative
\(V_{n+1}\).  The differentiated equation separates one time derivative from
two spatial derivatives, so the natural finite unit is an adjacent Sobolev
pair centred on the space in which the trace is sought.

\section{From adjacent rows to finite rectangles}

Fix a temporal row \(n\) and a spatial order \(m\), and write
\begin{equation*}
 v:=V_n,
 \qquad
 v_t=V_{n+1}.
\end{equation*}
The finite rectangle will place these two functions in the spaces
\begin{equation}
 v\in L^2(0,T_*;H^{m+1}),
 \qquad
 v_t\in L^2(0,T_*;H^{m-1}).
 \label{eq:adjacent-pair-prototype}
\end{equation}
The offset is deliberate.  The two spaces sit on opposite sides of
\(H^m\):
\begin{equation*}
 H^{m+1}\hookrightarrow H^m\hookrightarrow H^{m-1}.
\end{equation*}
At \(m=0\), this is the familiar triple
\begin{equation*}
 H^1\hookrightarrow L^2\hookrightarrow H^{-1}.
\end{equation*}
At higher \(m\), the same geometry is shifted up the Sobolev scale.

The middle space appears from an exact Fourier pairing.  For
\(f\in H^{m-1}\) and \(g\in H^{m+1}\), define
\begin{equation*}
 \langle f,g\rangle_m
 :=\int_{\R^3}(1+|\xi|^2)^m
   \widehat f(\xi)\cdot\overline{\widehat g(\xi)}\,d\xi.
\end{equation*}
Split the weight at the middle exponent:
\begin{equation*}
 (1+|\xi|^2)^m
 =(1+|\xi|^2)^{(m-1)/2}
  (1+|\xi|^2)^{(m+1)/2}.
\end{equation*}
Cauchy--Schwarz gives
\begin{equation}
 |\langle f,g\rangle_m|
 \le\norm{f}_{H^{m-1}}\norm{g}_{H^{m+1}}.
 \label{eq:adjacent-middle-pairing}
\end{equation}
For a smooth curve \(v\), differentiation of the middle norm yields
\begin{equation*}
 \frac d{dt}\frac12\norm{v(t)}_{H^m}^2
 =\operatorname{Re}\langle v_t(t),v(t)\rangle_m.
\end{equation*}
Integrating from \(r\) to \(t\) and applying
\eqref{eq:adjacent-middle-pairing} gives
\begin{align}
 \left|
 \norm{v(t)}_{H^m}^2-\norm{v(r)}_{H^m}^2
 \right|
 &\le2\int_r^t
 \norm{v_s(s)}_{H^{m-1}}
 \norm{v(s)}_{H^{m+1}}\,ds
 \notag\\
 &\le2
 \norm{v_t}_{L^2(r,t;H^{m-1})}
 \norm{v}_{L^2(r,t;H^{m+1})}.
 \label{eq:adjacent-norm-variation}
\end{align}
The upper spatial row and the next temporal row meet exactly at the
\(H^m\)-energy.  This is why the two Sobolev levels differ by two spatial
derivatives.

The same offset agrees with the parabolic order of the equation.  In
\eqref{eq:velocity-time-jet}, one new time derivative is balanced against
\(\Delta V_n\), which contains two new spatial derivatives.  The Hilbert
pairing explains how that parabolic separation produces a trace in the middle
space.

This is the Lions--Magenes trace mechanism in the shifted Sobolev triple; the
proof is included for the exact indices used here\bibref{LionsMagenes}.

\begin{lemma}[Adjacent-row trace]
\label{lem:part-ii-adjacent-trace}
Let \(m\in\R\) and let \(I=(0,T)\) with \(T<\infty\).  Suppose
\begin{equation*}
 v\in L^2(I;H^{m+1}),
 \qquad
 v_t\in L^2(I;H^{m-1})
\end{equation*}
in the distributional sense.  Then \(v\) has a unique representative in
\begin{equation*}
 C([0,T];H^m).
\end{equation*}
For that representative, estimate
\eqref{eq:adjacent-norm-variation} holds for every
\(0\le r\le t\le T\), and
\begin{equation}
 \begin{aligned}
 \norm{v}_{C([0,T];H^m)}^2
 &\le \frac1T\norm{v}_{L^2(I;H^m)}^2
 \\
 &\quad+2
 \norm{v_t}_{L^2(I;H^{m-1})}
 \norm{v}_{L^2(I;H^{m+1})}.
 \end{aligned}
 \label{eq:adjacent-trace-quantitative}
\end{equation}
\end{lemma}

\begin{proof}
Let \(S_R\) be a smooth Fourier cutoff to \(|\xi|\le2R\) that equals one on
\(|\xi|\le R\), and set
\begin{equation*}
 v_R:=S_Rv.
\end{equation*}
The cutoff commutes with \(\partial_t\).  At fixed \(R\), it raises spatial
regularity by any finite amount, so
\begin{equation*}
 v_R\in W^{1,2}(I;H^{m+1})
 \hookrightarrow C([0,T];H^{m+1}).
\end{equation*}
The smooth-row energy identity holds for \(v_R\).

Choose a Lebesgue time \(t_0\in I\) at which
\(v(t_0)\in H^{m+1}\).  Such times have full measure.  For two cutoffs
\(R,S\), apply the energy identity to
\(w=v_R-v_S\).  For every \(t\in[0,T]\),
\begin{align*}
 \norm{w(t)}_{H^m}^2
 &\le\norm{w(t_0)}_{H^m}^2
 \\
 &\quad+2
 \norm{w_t}_{L^2(I;H^{m-1})}
 \norm{w}_{L^2(I;H^{m+1})}.
\end{align*}
As \(R,S\to\infty\), the first term tends to zero because
\(v(t_0)\in H^{m+1}\).  The two tail norms also tend to zero because Fourier
cutoffs converge strongly in the given Bochner spaces.  Hence
\((v_R)_R\) is Cauchy in \(C([0,T];H^m)\).  Its uniform limit
\(\widetilde v\) is continuous into \(H^m\).

On the other hand, \(v_R\to v\) in
\(L^2(I;H^{m+1})\), hence in \(L^2(I;H^m)\).  The continuous limit
\(\widetilde v\) agrees with \(v\) almost everywhere and is its required
representative.  Passing to the limit in the cutoff energy identity proves
\eqref{eq:adjacent-norm-variation} for \(\widetilde v\).  Two continuous
representatives that agree almost everywhere agree everywhere, which gives
uniqueness.

The continuous representative belongs to \(L^2(I;H^m)\), so some
\(r\in[0,T]\) satisfies
\[
 \norm{v(r)}_{H^m}^2
 \le \frac1T\norm{v}_{L^2(I;H^m)}^2.
\]
Apply \eqref{eq:adjacent-norm-variation} on the interval between this \(r\)
and an arbitrary \(t\in[0,T]\), then take the supremum over \(t\).  This gives
\eqref{eq:adjacent-trace-quantitative}.
\end{proof}

The lemma controls the whole vector in the middle space, while the displayed
estimate controls the change of its squared norm.  The cutoff construction
joins those two facts in one continuous representative.  In particular, if
the adjacent pair holds on \((0,T_*)\), then applying the lemma with
\(T=T_*\) produces one value \(v(T_*)\in H^m\), and every sequence
\(t_k\uparrow T_*\) converges to that same value.  The topology
selects one compatible terminal reading and resolves the failure called Jump
in Part~I.

To use this at every depth, the adjacent pair must be retained for each
temporal row and each middle Sobolev order, together with the pressure produced
by those rows.  The smallest continuation-grade example is the pair for \(u\)
centred on \(H^2\).

\subsection{The first concrete rectangles}

Choose temporal cutoff \(N=0\) and middle spatial order \(M=2\).  The
resulting rectangle has velocity size
\begin{equation}
 \mathcal N_{0,2}(u;T)
 :=\norm{u}_{L^2(0,T;H^3)}^2
   +\norm{u_t}_{L^2(0,T;H^1)}^2.
 \label{eq:first-concrete-rectangle}
\end{equation}
The norm \eqref{eq:first-concrete-rectangle} is the adjacent pair
\eqref{eq:adjacent-pair-prototype} placed inside
the mixed jet.  Each coordinate pays a different part of the terminal
argument.

First, the two coordinates control the total variation of the critical
height.  The kinematic identity at \(n=1\) gives
\begin{equation*}
 H'(t)=\operatorname{Re}
 \pair{\Lambda^{1/2}u_t(t)}{\Lambda^{1/2}u(t)}.
\end{equation*}
Since
\(H^1\hookrightarrow\dot H^{1/2}\) and
\(H^3\hookrightarrow\dot H^{1/2}\), Cauchy--Schwarz in time gives
\begin{align*}
 \int_0^{T_*}|H'(t)|\,dt
 &\le
 \norm{u_t}_{L^2(0,T_*;\dot H^{1/2})}
 \norm{u}_{L^2(0,T_*;\dot H^{1/2})}
 \\
 &\le
 C\norm{u_t}_{L^2(0,T_*;H^1)}
   \norm{u}_{L^2(0,T_*;H^3)}.
\end{align*}
Finiteness of the whole-terminal rectangle would make the right side finite.
Then
\begin{equation*}
 H(t)\le H(0)+\int_0^{T_*}|H'(s)|\,ds
\end{equation*}
would rule out the terminal escape in
\eqref{eq:critical-height-limsup}.

The same two coordinates also pay the signed passage bill derived earlier.
The upper row gives
\begin{equation*}
 \int_0^{T_*}E_{3/2}(t)\,dt
 \le \frac12\norm{u}_{L^2(0,T_*;H^3)}^2<\infty,
\end{equation*}
and the critical balance \eqref{eq:H1} then gives
\begin{equation*}
 \mathcal P_{1/2}=H'+2\nu E_{3/2}\in L^1(0,T_*).
\end{equation*}
For the dyadic first-passage intervals,
\eqref{eq:dyadic-passage-work} requires
\[
 L_j
 \le \int_{t_j}^{t_{j+1}}\mathcal P_{1/2}(s)\,ds
 \le \int_{t_j}^{T_*}|\mathcal P_{1/2}(s)|\,ds.
\]
The last tail tends to zero because \(t_j\uparrow T_*\), while
\(L_j\to\infty\).  Thus the first rectangle closes not only the scalar
variation estimate but the exact VPI work requirement that arose from Zeno
compression.

Second, the upper coordinate produces the continuation-grade Lipschitz
coefficient.  Sobolev embedding gives
\begin{equation*}
 H^3(\R^3)\hookrightarrow W^{1,\infty}(\R^3),
\end{equation*}
and Cauchy--Schwarz on the finite interval gives
\begin{equation*}
 \int_0^{T_*}\norm{\nabla u(t)}_\infty\,dt
 \le C\sqrt{T_*}\,
 \norm{u}_{L^2(0,T_*;H^3)}.
\end{equation*}
This is the integrable coefficient required by the higher Sobolev continuation
estimate in Part~I.

Third, the two rows surround the middle space \(H^2\).  Applying
Lemma~\ref{lem:part-ii-adjacent-trace} with \(m=2\) gives
\begin{equation*}
 u\in C([0,T_*];H^2).
\end{equation*}
The first rectangle consequently controls critical variation, produces the
Lipschitz coefficient, and gives the original velocity curve a strong terminal
trace.  One finite pair produces all three consequences together.

Raise the temporal cutoff to \(N=1\) while keeping \(M=2\).  The next
rectangle has size
\begin{equation*}
\begin{aligned}
 \mathcal N_{1,2}(u;T)
 :={}&\norm{u}_{L^2(0,T;H^3)}^2
     +\norm{u_t}_{L^2(0,T;H^3)}^2
 \\
 &+\norm{u_t}_{L^2(0,T;H^1)}^2
     +\norm{u_{tt}}_{L^2(0,T;H^1)}^2.
\end{aligned}
\end{equation*}
The shared coordinate \(u_t\) appears in both spatial rows.  In the first
adjacent pair it is the derivative of \(u\); in the second it is the evolving
field whose derivative is \(u_{tt}\).  Both appearances refer to the same
distributional derivative generated by the equation.  Raising \(N\) extends
this overlap through higher temporal rows.  Raising \(M\) places the same
rows around a stronger middle trace space.

For arbitrary finite \(N,M\in\Nzero\), define
\begin{equation}
\begin{aligned}
 \mathcal N_{N,M}(u;T)
 :={}&\sum_{n=0}^N
 \norm{V_n}_{L^2(0,T;H^{M+1})}^2
 \\
 &+\sum_{n=0}^N
 \norm{V_{n+1}}_{L^2(0,T;H^{M-1})}^2.
\end{aligned}
\label{eq:rectangle-norm}
\end{equation}
When \(M=0\), the lower space is \(H^{-1}\).  For every \(M\), the index
names the middle trace space \(H^M\); pointwise derivative counts are a
different notion.  Each cutoff pair carries its own constant, and the family
is combined coordinatewise.

A bound uniform in the time cutoff \(T\), for this fixed pair \((N,M)\), has
two immediate consequences.  It supplies every adjacent pair in
\eqref{eq:rectangle-norm}, so Lemma~\ref{lem:part-ii-adjacent-trace} gives
\begin{equation*}
 V_n\in C([0,T_*];H^M),
 \qquad
 \sup_{0\le t\le T_*}\norm{V_n(t)}_{H^M}<\infty
 \quad(0\le n\le N).
\end{equation*}
At the chosen cutoff, every traced row remains bounded in its middle Sobolev
topology, which is the finite-depth control demanded by Blown.  The family
ranges over every finite \(M\), so every finite spatial readout is reached by
a larger rectangle.  The same continuous representative forces each row to
approach one value in \(H^M\), resolving Jump at that cutoff.  The rectangle
keeps both consequences attached to the same derivatives at every depth.

\subsubsection{The pressure-complete finite rectangle space}

Let \(I=(0,T)\), where \(0<T\le T_*\).  The velocity space belonging to
the cutoff pair \((N,M)\) is
\begin{equation*}
 \mathfrak V_{N,M}(I)
 :=\left\{
 (v_0,\ldots,v_{N+1}):
 \begin{array}{ll}
 v_n\in L^2(I;H^{M+1}_\sigma),&0\le n\le N,\\
 v_{n+1}\in L^2(I;H^{M-1}_\sigma),&0\le n\le N,\\
 \partial_tv_n=v_{n+1}\text{ in }\mathcal D'(I\times\R^3),&0\le n\le N
 \end{array}
 \right\}.
\end{equation*}
The norm is the finite sum displayed in \eqref{eq:rectangle-norm}.  Each
\(v_n\), \(0\le n\le N\), has an \(H^M\)-continuous representative by the
adjacent-row lemma and has a well-defined trace at \(t=0\).

The pressure-complete rectangle space
\(\mathscr R_{N,M}(u_0;I)\) contains the tuples
\begin{equation*}
 (v_0,\ldots,v_N;v_{N+1};q_0,\ldots,q_N)
\end{equation*}
in
\begin{equation*}
 \mathfrak V_{N,M}(I)
 \times[L^1(I;H^M)]^{N+1}
\end{equation*}
that satisfy four finite conditions:
\begin{enumerate}
 \item \(\partial_tv_n=v_{n+1}\) and \(\nabla\cdot v_n=0\) for
       \(0\le n\le N\);
 \item the pressure recurrence \eqref{eq:pressure-time-jet} holds with
       \(V_a=v_a\) and \(Q_n=q_n\);
 \item the momentum recurrence \eqref{eq:velocity-time-jet} holds in
       \(\mathcal D'(I\times\R^3)\);
 \item the traces of \(v_0,\ldots,v_N\) at the initial time equal the
       datum-generated values from
       \eqref{eq:initial-pressure}--\eqref{eq:initial-velocity}.
\end{enumerate}
At this cutoff, \(v_{N+1}\) has only its lower-space membership.  Its
\(H^M\) trace appears when the temporal cutoff is raised and the same function
is retained as an upper row.  The pressure coordinate lies in
\(L^1_tH^M_x\) and inherits its value from the velocity rows through the
Riesz formula.

For the maximal classical solution, the finite rectangle realized on
\(I_T=(0,T)\) is
\begin{equation*}
 \cR_{N,M}[u_0;T]
 :=(V_0,\ldots,V_N;V_{N+1};Q_0,\ldots,Q_N)|_{I_T}
 \in\mathscr R_{N,M}(u_0;I_T).
\end{equation*}
The notation \(u_0\) records the origin of every entry.  The rectangle is a finite
restriction of the jet generated by \(u_0\) and realized by the maximal
solution.

\subsubsection{The velocity rows determine the pressure topology}

The product estimates explain why the pressure coordinate belongs to
\(L^1(I;H^M)\).  When \(M=0\), interpolation between
\(H^1\hookrightarrow L^2\) and \(H^1\hookrightarrow L^6\) gives
\begin{equation*}
 \norm f_4
 \le\norm f_2^{1/4}\norm f_6^{3/4}
 \le C\norm f_{H^1}.
\end{equation*}
Consequently,
\begin{equation*}
 \norm{fg}_2
 \le\norm f_4\norm g_4
 \le C\norm f_{H^1}\norm g_{H^1}.
\end{equation*}
When \(M\ge1\), the order \(M+1\) is greater than \(3/2\), and
\(H^{M+1}(\R^3)\) is a Banach algebra.  In both cases,
\begin{equation}
 \norm{fg}_{H^M}
 \le C_M
 \norm f_{H^{M+1}}\norm g_{H^{M+1}}.
 \label{eq:rectangle-product}
\end{equation}
The Riesz transforms are bounded on \(H^M\).  Applying
\eqref{eq:rectangle-product} to every finite product in
\eqref{eq:pressure-time-jet}, then using Cauchy--Schwarz in time, gives
\begin{equation}
 \norm{Q_n}_{L^1(J;H^M)}
 \le C_{n,M}\sum_{a=0}^n
 \norm{V_a}_{L^2(J;H^{M+1})}
 \norm{V_{n-a}}_{L^2(J;H^{M+1})}
 \label{eq:rectangle-pressure-topology}
\end{equation}
for every subinterval \(J\subset(0,T_*)\) on which the displayed velocity
norms are finite.

The differentiated momentum equation also has a common topology.  Since every
\(V_a\) is divergence-free,
\begin{equation*}
 (V_a\cdot\nabla)V_b
 =\nabla\cdot(V_a\otimes V_b).
\end{equation*}
At \(M=0\), the \(H^1\hookrightarrow L^4\) estimate puts
\(V_a\otimes V_b\) in \(L^1_tL^2_x\), so its divergence lies in
\(L^1_tH^{-1}_x\).  At \(M\ge1\), the algebra estimate puts the tensor
product in \(L^1_tH^{M+1}_x\), and its divergence belongs to
\(L^1_tH^M_x\), hence to \(L^1_tH^{M-1}_x\).  The pressure gradient belongs
to \(L^1_tH^{M-1}_x\) by
\eqref{eq:rectangle-pressure-topology}.  Finally,
\begin{equation*}
 \Delta V_n\in L^2(I;H^{M-1})
 \hookrightarrow L^1(I;H^{M-1})
\end{equation*}
because \(I\) has finite length.  The complete temporal recurrence
holds in \(L^1(I;H^{M-1})\).

The upper velocity rows fix the pressure size and gauge.  Its gradient meets
the lower velocity topology, and the same finite rectangle retains the
incompressibility equation that generated it.

The finite rectangle now retains the complete pressure-coupled response needed
at the endpoint.  Its upper velocity rows determine every normalized pressure
row through \eqref{eq:pressure-time-jet}; estimate
\eqref{eq:rectangle-pressure-topology} places that pressure in the common
topology; and every term of \eqref{eq:velocity-time-jet} remains meaningful in
\(L^1(I;H^{M-1})\).  Whole-terminal bounds and compatible traces can then carry
the velocity, pressure, and temporal response to one endpoint, where the
limiting recurrence rules out Dead.

\subsection{Compatibility under restriction}

A larger rectangle can belong to the same all-orders family only if lowering
either cutoff retains the functions already present.  Suppose \(N'\le N\)
and \(M'\le M\).  Deleting rows above \(N'\) leaves
\(V_0,\ldots,V_{N'+1}\) and \(Q_0,\ldots,Q_{N'}\) unchanged, while the
embeddings
\[
 H^{M+1}\hookrightarrow H^{M'+1},\qquad
 H^{M-1}\hookrightarrow H^{M'-1},\qquad
 H^M\hookrightarrow H^{M'}
\]
view those same functions at the lower spatial order.

For example, the \((2,3)\) rectangle contains
\(u,u_t,u_{tt},u_{ttt}\) and \(p,p_t,p_{tt}\).  Lowering the cutoffs to
\((1,2)\) discards \(u_{ttt}\) and \(p_{tt}\), retains
\(u,u_t,u_{tt}\) and \(p,p_t\), and regards every retained function in the
weaker Sobolev topology.  The retained \(u_t\) is the same distributional
derivative in both rectangles.

This equality on every shared derivative is compatibility.  Once the bounds
hold on the common terminal interval, these restriction identities assemble
the finite rectangles into one all-orders object.  Each coordinate retains
its own constant indexed by finite derivative depth.

The same equality later forces agreement of the terminal traces.  If two
rectangles contain the row \(V_n\), both use the same preterminal curve.  When
their adjacent estimates produce continuous representatives in a shared lower
Sobolev topology, those representatives agree for every \(t<T_*\) and hence at
\(T_*\).  Compatibility is thus the reason a derivative cannot acquire one
endpoint value in a strong rectangle and another value in a weaker one.

\subsection{The terminal quantifier}

Compatibility fixes the identity of the coordinates.  Terminal control adds a
second requirement: each fixed finite rectangle must remain bounded as its
time interval approaches \(T_*\).

Fix \(T<T_*\), a temporal row \(n\), and a spatial order \(q\).  The maximal
solution is classical on the compact interval \([0,T]\), so
\begin{equation*}
 \sup_{0\le t\le T}\norm{V_n(t)}_{H^q}<\infty.
\end{equation*}
It follows that
\begin{equation*}
 \norm{V_n}_{L^2(0,T;H^q)}^2
 \le T\sup_{0\le t\le T}\norm{V_n(t)}_{H^q}^2
 <\infty.
\end{equation*}
Every sum defining \(\mathcal N_{N,M}(u;T)\) contains only finitely many such
terms.  Hence local classicality gives
\begin{equation}
 \forall T<T_*\ \forall N,M<\infty,
 \qquad
 \mathcal N_{N,M}(u;T)<\infty.
 \label{eq:strict-subinterval-rectangles}
\end{equation}
The bound obtained from compactness may depend on \(T\).  All of the proposed
singular behavior can be carried by that dependence.

The scalar function
\begin{equation*}
 f(t)=(T_*-t)^{-1/2}
\end{equation*}
isolates the point.  For each fixed \(T<T_*\),
\begin{equation*}
 \int_0^T|f(t)|^2\,dt
 =\log\frac{T_*}{T_*-T}<\infty.
\end{equation*}
On the complete interval,
\begin{equation*}
 \int_0^{T_*}|f(t)|^2\,dt
 =\int_0^{T_*}\frac{dt}{T_*-t}
 =\infty.
\end{equation*}
The example is a scalar model of the quantifier failure.  Finiteness on every
strict subinterval coexists here with divergence on their union.

Local classicality says that \(\mathcal N_{N,M}(u;T)\) is finite after a cutoff
\(T<T_*\) has been fixed; its finite value may diverge as the cutoff approaches
the endpoint.  The required statement begins with the rectangle
directly on \(I_*:=(0,T_*)\).  Write \(\mathcal N_{N,M}(u;I_*)\) for the finite sum in
\eqref{eq:rectangle-norm} with every time norm taken on \(I_*\).  The
whole-terminal statement is
\begin{equation}
 \mathcal N_{N,M}(u;I_*)<\infty
 \qquad\text{for every finite }N,M.
 \label{eq:whole-terminal-quantifier}
\end{equation}
This direct datum-generated whole-terminal finiteness is the established theorem that opens
Part~III, supplies the rectangle family, completes its pressure rows, and
assembles its shared derivatives.
